\section{Discussion}
\label{sec:discussion}

The results in Section~\ref{sec:results} are most usefully read as
three separable themes: the magnitude--direction asymmetry of the
news response, the methodological consequences of pre-event drift,
and the limits of the present design relative to extensions we do
not pursue.

\subsection{Magnitude is easier to predict than direction}

The result with the most uniform statistical signal across cells is
that unscheduled news events are associated with substantially
elevated intraday magnitude (Section~\ref{ssec:results-volatility}),
while the directional content of the same headlines is small and
too weak to treat as a standalone signal
(Section~\ref{ssec:results-direction}). The two results are not in
tension: they describe complementary properties of the same
underlying phenomenon. A breaking headline produces price
discovery---the market revises its mid-quote in response to new
information---but the sign of the revision depends on context, prior
positioning, and the fraction of the news that was already
anticipated. Magnitude captures the price-discovery event itself;
direction asks the harder question of which side of the new
mid-quote the prior expectation was on.

From an applied perspective, news of this kind is more useful as a
volatility filter than as a directional signal: volatility-aware
position sizing, stop-loss design that allows for post-news range
expansion, and execution rules that avoid mean-reversion trades
immediately after a headline are natural operational uses of the
magnitude finding. The directional content of the LLM sentiment, by
contrast, is best treated as one weak predictor among many rather
than as a standalone trigger. The pattern is consistent with
\citet{antweiler2004talk}, who find that message-board posting
activity (a volume proxy) is more strongly associated with realised
volatility than its bullishness content is with returns, and with
the \citet{heston2017news} observation that news has faster effects
on volatility than on directional returns.

The hit-rate advantage of the LLM-derived sentiment over the
dictionary baseline is about 1--6 percentage points on the primary
windows, and the LLM is competitive with FinBERT, beating it in
five of six primary cells (Section~\ref{ssec:results-direction}).
This is not a dramatic improvement on the headline metric. The qualitative
advantage is more substantial: the LLM produces different sentiment
for USD and for NDX on $73.5\%$ of events, a property that captures
the cross-asset asymmetry of risk-on/risk-off events that a
single-polarity classifier cannot express. We recommend that
subsequent work in this area report both metrics jointly, so that
the cross-asset coherence of LLM labels with the known
macro-finance taxonomy can be evaluated alongside the per-asset hit
rate.

\subsection{Pre-event drift: latency, not necessarily front-running}

The pre-event drift result (Section~\ref{ssec:results-timing}, H8)
is the methodologically most consequential finding of the analysis.
The absolute return in the fifteen minutes preceding a Discord
headline is elevated above the matched baseline at
$q = 7.5\times10^{-47}$ for EUR/USD and
$q = 2.2\times10^{-40}$ for NDX, with a magnitude comparable to
the post-event drift on the same window.

A naive interpretation would invoke insider trading or systematic
front-running. We caution against that reading. The Discord
timestamp on which our event windows are aligned is downstream of
the underlying primary source---typically a wire service, an agency
feed, or an official release calendar---and the variable latency
between the primary source and the Discord publication is, in our
reading, the most plausible explanation. Without direct access to
the primary-source timestamp for each event we cannot decompose the
pre-event drift into ``feed latency'' and ``true pre-disclosure
activity''; we expect that the former dominates for most events.

The implication is the same regardless of the decomposition: the
standard event-study assumption that the public publication time
coincides with the informational event time is materially violated
in this setting, and the apparent reaction in any post-event window
includes a component that has, in fact, already taken place.
Subsequent work on online-news event studies should either obtain
primary-source timestamps or report the pre-event drift as a
diagnostic alongside the post-event response.

\subsection{What we do not test}

We do not test whether more sophisticated LLM prompting
(chain-of-thought, multi-prompt ensembles), more recent LLM models,
or fine-tuning on a hand-labelled subsample would improve the
directional results beyond the present small single-digit
percentage-point margins. We do not
extend the asset universe beyond \eurusd{} spot and the
\ndx{} CFD, nor do we test individual equities where firm-specific
news effects would dominate. We do not decompose the magnitude
response into a re-pricing component and a liquidity-withdrawal
component, which would require an order-flow or quote-imbalance
dataset we do not have. And we do not test cross-feed replication on
a second public news service, which would quantify the
feed-specific bias of the FinancialJuice gold filter. We view these
extensions, particularly in combination with primary-source
timestamps that would allow the pre-event drift to be decomposed
cleanly, as the most promising directions for follow-up work.
